\chapter{Texto \texorpdfstring{$\to$}{->} Dados: JSON, Esquemas e ETL Orientado a Prompts}
\label{ch:texto-dados}

\ficcao O apoio ao cliente da Lavandaria recebe «A máquina ficou com a minha meia. E a fatura ainda não chegou.» Para responder, podes ler a mensagem inteira. Para contar os assuntos do mês, precisas de identificar os dois pedidos e guardar os campos que permitem somá-los.

\section{Escolher os campos antes de pedir a extração}
Começa por escrever à mão a saída de dois ou três casos. Que campo pode ficar vazio? A categoria é escolhida de uma lista ou pode ser texto livre? Um esquema descreve a estrutura que aceitamos, incluindo os tipos e os valores admitidos.

ETL significa extrair, transformar e carregar dados. Num processo de apoio, podes ler a mensagem, convertê-la num objeto com campos e carregar os registos validados num sistema de trabalho. Um modelo pode ajudar na extração quando a informação está em linguagem livre. Para converter uma data já estruturada ou somar casos, usa a regra ou a consulta que consegues testar.

Exemplo: gerir pedidos de suporte por e-mail. Em vez de “resumir”, defines:

\begin{itemize}[leftmargin=*]
\item \texttt{id\_cliente} (string),
\item \texttt{tipo\_problema} (enum: “pagamento”, “login”, “bug”, “outro”),
\item \texttt{urgencia} (enum: “baixa”, “media”, “alta”),
\item \texttt{descricao\_curta} (string),
\item \texttt{precisa\_resposta\_humana} (boolean),
\item \texttt{sentimento} (enum: “negativo”, “neutro”, “positivo”).
\end{itemize}

Pede apenas os campos necessários à decisão. Se extraíres um identificador de cliente da mensagem, confirma-o com os registos autorizados antes de o usar para ligar uma reclamação a uma conta.

\begin{promptlabbox}{Definir schema antes de pedir extração}
\textbf{Objetivo e contexto:} Escolher campos de extração para três a cinco textos fictícios e criar referências manuais.

Escolhe um tipo de texto que tenhas em massa (e-mails, comentários, notas de reuniões).

\medskip

\textbf{Passo 1} — Decide 5–8 campos que seriam úteis para análise. Escreve-os como numa tabela: nome, tipo, descrição.

\textbf{Passo 2} — Só depois disso, escreve um prompt do tipo:

\begin{lstlisting}
A partir do texto abaixo, extrai os seguintes campos:

id_cliente: ...

tipo_problema: um entre ["pagamento", "login", "bug", "outro"]

urgencia: um entre ["baixa", "media", "alta"]

descricao_curta: ...

precisa_resposta_humana: true ou false

sentimento: um entre ["negativo", "neutro", "positivo"]

Devolve APENAS um objeto JSON válido.
\end{lstlisting}

Aplica aos mesmos três a cinco textos fictícios ou autorizados e verifica se o modelo respeita o esquema.

\textbf{Resultado esperado e validação:} Faz parsing e confere campos, tipos, enumerações e informação ausente. Compara valores com os textos, sem concluir que JSON válido implica verdade.

\textbf{Falhas possíveis:} Preencher campos não presentes, usar identificadores reais ou avaliar apenas se há chavetas.
\end{promptlabbox}

\section{JSON como contrato de saída}

JSON representa dados como texto. Um objeto junta pares de nome e valor entre chavetas. Uma lista agrupa valores entre parênteses retos. O texto fica entre aspas duplas. Valores como \texttt{true}, \texttt{false} e \texttt{null} têm significado próprio no formato~\citep{rfc8259}.

Um exemplo de saída ajuda a mostrar a forma. Ao receber o resultado, o programa tenta lê-lo, verifica os campos e aplica as regras do domínio. Se o modelo escrever «Aqui está» antes do objeto, o ficheiro completo pode falhar no primeiro passo.

\begin{promptlabbox}{Forçar JSON limpo}
\textbf{Objetivo e contexto:} Separar pedido de JSON e validação. Usa as mesmas mensagens fictícias e um esquema explícito.

Tenta um prompt deste tipo:

\medskip

\begin{lstlisting}
Tarefa: extrair informação de um texto de suporte.

Lê o texto do cliente.

Preenche um objeto JSON com o seguinte formato:

{
"tipo_problema": "bug",
"urgencia": "alta",
"descricao_curta": "string",
"precisa_resposta_humana": true
}

IMPORTANTE:

Devolve APENAS o JSON, sem texto antes ou depois.

Garante que o JSON é válido (aspas, vírgulas, etc.).

Texto:
[cola aqui o e-mail do cliente]
\end{lstlisting}

O bloco é uma instância concreta; as opções admitidas devem ser definidas num esquema separado. Faz parsing local com código, sem enviar dados pessoais para um validador público. Regista também os casos de recusa, interrupção e resposta vazia.

\textbf{Resultado esperado e validação:} Executa parsing local e valida o esquema, incluindo campos extra e ausência de obrigatórios. Regista respostas vazias, recusas e interrupções.

\textbf{Falhas possíveis:} Usar \texttt{True} ou \texttt{False}, com maiúscula, em JSON, enviar dados privados para validadores públicos ou ignorar validade semântica.
\end{promptlabbox}

\section{Lidar com incerteza e campos em falta}

Texto pode ser irregular e ambíguo. Nem sempre o e-mail vai dizer claramente a urgência ou o tipo de problema. Se o esquema exigir uma categoria mesmo quando falta informação, uma saída pode cumprir a forma com um palpite.

Em vez disso, podes:

\begin{itemize}[leftmargin=*]
\item permitir valores \texttt{null},
\item introduzir uma opção \texttt{"desconhecido"},
\item pedir um excerto de evidência e permitir encaminhamento humano.
\end{itemize}

Um esquema pode exigir o campo \texttt{evidencia} e admitir \texttt{null} para indicar que falta suporte. Se o campo estiver ausente, temos outro estado. Combina esses significados antes de automatizar a decisão. O objeto seguinte mostra uma instância que o esquema teria de admitir.

\begin{lstlisting}
{
"tipo_problema": "desconhecido",
"urgencia": "desconhecida",
"descricao_curta": "string",
"precisa_resposta_humana": true,
"evidencia": null
}
\end{lstlisting}

No prompt, explicas:

\begin{lstlisting}
Se não conseguires inferir um campo com segurança,
usa "desconhecido" ou null; identifica o dado em falta.
Nao inventes probabilidades de acerto.
\end{lstlisting}

\begin{warningboxenv}{Não obrigues a IA a adivinhar só para encher campos}
Schemas sem opção de “desconhecido” são convites à alucinação. Em dados de produção, prefiro um campo vazio honesto a um valor inventado com ar respeitável.
\end{warningboxenv}

\section{Validação e pós-processamento de dados extraídos}

Mesmo com prompts bons, vai haver erros. Então assume que:

\begin{itemize}[leftmargin=*]
\item o JSON precisa de validação,
\item campos enumerados podem vir com valores fora da lista,
\item números podem vir como strings, etc.
\end{itemize}

Padrão saudável de pós-processamento:

\begin{enumerate}[leftmargin=*]
\item Faz parsing do JSON com tolerância mínima (se falhar, regista erro e, se fizer sentido, tenta uma segunda passagem de correção).
\item Valida cada campo:
\begin{itemize}[leftmargin=*]
\item se o valor está fora da enumeração, rejeita ou encaminha segundo uma regra explícita,
\item se um número vem como string, só converte quando o contrato documenta essa normalização,
\item se está completamente fora do esperado, regista para análise manual.
\end{itemize}
\item Mede estatísticas de erro:
\begin{itemize}[leftmargin=*]
\item percentagem de registos com erros graves,
\item tipos mais comuns de erro,
\item campos mais problemáticos.
\end{itemize}
\end{enumerate}

\begin{checklistbox}{Pipeline robusto de texto \texorpdfstring{$\to$}{->} dados}
Para cada fluxo ETL com LLM:

\begin{itemize}[leftmargin=*]
\item \textbf{Schema explícito}: campos, tipos, enums, possibilidade de “desconhecido”.
\item \textbf{Prompt claro}: exemplos de JSON, instruções para lidar com incerteza.
\item \textbf{Validação automática}: parsing, tipos, enums, intervalos.
\item \textbf{Logs de erro}: guardas casos falhados para análise e melhoria de prompts.
\item \textbf{Monitorização}: medes qualidade ao longo do tempo (taxa de falhas, campos mais problemáticos).
\end{itemize}
\end{checklistbox}

\section{Redução de ambiguidade com instruções positivas}

Ao desenhar prompts de extração, evita instruções vagas tipo “identifica o tipo de problema”. Em vez disso, usa instruções positivas e listas fechadas:

\begin{lstlisting}
tipo_problema deve ser UMA das seguintes opções:

"pagamento" se o problema envolver faturas, cartões,
cobranças, débitos, reembolsos.

"login" se envolver palavras-passe, autenticação,
contas bloqueadas.

"bug" se o utilizador descrever falhas técnicas
na aplicação, erros, crashes.

"outro" para todos os restantes casos que não
se encaixem nas anteriores.
\end{lstlisting}

Isto:

\begin{itemize}[leftmargin=*]
\item reduz a criatividade indesejada,
\item torna o comportamento mais previsível,
\item facilita explicar ao negócio o que significa cada valor.
\end{itemize}

\begin{keyideabox}{Especificar é cansativo, mas paga-se em estabilidade}
Quanto mais preguiça tiveres a escrever critérios no prompt, mais trabalho vais ter depois a limpar os dados. Escolhe onde queres gastar energia.
\end{keyideabox}

\section{Validade sintática, estrutural e semântica}
\label{sec:validacao-json}
O primeiro passo é descodificar o texto. O RFC~8259 exclui \texttt{NaN} e recomenda nomes únicos em cada objeto. Por omissão, o módulo \texttt{json} do Python aceita \texttt{NaN} e conserva o último valor de um nome repetido~\citep{rfc8259,pythonjson312}. Verificámos este exemplo em Python 3.12.14.

\begin{lstlisting}
>>> import json
>>> json.loads('{"prazo": 10, "prazo": 5}')
{'prazo': 5}
\end{lstlisting}

Usa \texttt{object\_pairs\_hook} no descodificador para rejeitar nomes repetidos quando os dois pares ainda existem. Com \texttt{parse\_constant}, podes rejeitar as constantes numéricas excluídas pelo RFC. Depois de o texto passar esse ponto, JSON Schema verifica os tipos e as restrições declaradas~\citep{jsonschema2020}. A aplicação ainda tem de verificar se o prazo corresponde à política e se a reclamação pertence ao cliente autorizado.

Algumas APIs oferecem saídas condicionadas por esquema~\citep{structured2026}. É uma capacidade útil, com subconjuntos de esquema e condições próprias do serviço. Não confundas «JSON mode», conformidade com esquema e correção factual. Trata recusas e gerações incompletas como estados explícitos.

No incidente do capítulo~\ref{ch:llmops}, o prazo antigo passa a verificação de forma. É a comparação com a política em vigor que impede a resposta. A seguir vamos trabalhar com textos que pedem ainda outra revisão, a da voz e da intenção de quem os escreve.
